% !TEX root = ../DoAn.tex
\documentclass[../DoAn.tex]{subfiles}
\begin{document}

\section{Công nghệ Front-end}
\label{section:3.1}
\subsection{Next.js và React trong xây dựng giao diện người dùng}
\label{subsection:3.1.1}

Frontend của hệ thống quản lý sinh viên được xây dựng bằng Next.js kết hợp với React \cite{react} nhằm tạo ra giao diện trực quan, hiện đại và dễ dàng mở rộng.

React là thư viện JavaScript hỗ trợ xây dựng giao diện người dùng theo mô hình component, cho phép tái sử dụng các thành phần giao diện và nâng cao hiệu quả phát triển. Trong khi đó, Next.js là một framework phát triển dựa trên React, cung cấp sẵn hệ thống định tuyến, tối ưu hiệu năng và hỗ trợ nhiều phương thức kết xuất như Server-Side Rendering (SSR), Static Site Generation (SSG) và Client-Side Rendering (CSR).

Việc sử dụng Next.js giúp tổ chức cấu trúc dự án rõ ràng, cải thiện tốc độ tải trang và nâng cao trải nghiệm người dùng.
\subsection{TypeScript trong phát triển ứng dụng}
\label{subsection:3.1.2}
Để tăng tính ổn định và giảm thiểu lỗi trong quá trình phát triển, hệ thống sử dụng TypeScript thay cho JavaScript thuần túy.

TypeScript bổ sung cơ chế kiểm tra kiểu dữ liệu tĩnh (static typing), giúp định nghĩa rõ ràng cấu trúc dữ liệu và phát hiện lỗi ngay trong quá trình lập trình. Điều này đặc biệt hữu ích đối với hệ thống quản lý sinh viên khi phải xử lý nhiều dữ liệu liên quan như thông tin sinh viên, lớp học, môn học và kết quả học tập.

Ngoài ra, TypeScript còn hỗ trợ tốt cho việc bảo trì, mở rộng và nâng cao năng suất phát triển thông qua các công cụ gợi ý mã nguồn (IntelliSense).
\subsection{Quản lý trạng thái với Zustand và TanStack Query}
\label{subsection:3.1.3}
Hệ thống sử dụng kết hợp Zustand và TanStack Query để quản lý dữ liệu phía client.

Zustand: Được sử dụng để quản lý các trạng thái cục bộ như thông tin người dùng đăng nhập, trạng thái giao diện và các thiết lập của ứng dụng.

TanStack Query (React Query): Được sử dụng để quản lý dữ liệu từ máy chủ, hỗ trợ gọi API, lưu bộ nhớ đệm (caching), đồng bộ dữ liệu và xử lý các trạng thái tải dữ liệu hoặc lỗi.

Sự kết hợp này giúp giảm số lượng yêu cầu gửi đến máy chủ, tối ưu hiệu năng và nâng cao trải nghiệm sử dụng của người dùng.

\subsection{Tailwind CSS và các thư viện hỗ trợ giao diện}
\label{subsection:3.1.4}
Hệ thống sử dụng Tailwind CSS để xây dựng giao diện theo phương pháp utility-first. Cách tiếp cận này giúp đảm bảo tính nhất quán trong thiết kế, tăng tốc độ phát triển và tối ưu dung lượng tệp CSS sau khi biên dịch
Bên cạnh đó, hệ thống còn sử dụng một số thư viện hỗ trợ như:

* React Hook Form: quản lý và xử lý dữ liệu biểu mẫu.

* Zod: kiểm tra và xác thực dữ liệu đầu vào.

* Recharts: trực quan hóa dữ liệu thống kê dưới dạng biểu đồ.

\section{Công nghệ Back-end và cơ sở dữ liệu}
\label{section:3.2}
\subsection{Express.js trong xây dựng hệ thống Backend}
\label{subsection:3.2.1}
Backend của hệ thống được xây dựng trên nền tảng Node.js \cite{nodejs} và sử dụng Express.js \cite{express} để phát triển các dịch vụ RESTful API.

Express.js \cite{express} là một framework tối giản, linh hoạt và dễ mở rộng, cho phép xây dựng các thành phần như Router, Middleware và Controller để xử lý các nghiệp vụ của hệ thống. Đồng thời, Express.js cũng hỗ trợ tốt việc kiểm soát luồng dữ liệu, xử lý lỗi và xác thực người dùng.
\subsection{PostgreSQL và Sequelize ORM}
\label{subsection:3.2.2}
Hệ thống sử dụng PostgreSQL \cite{postgresql} làm hệ quản trị cơ sở dữ liệu chính nhằm lưu trữ thông tin sinh viên, lớp học, môn học và các dữ liệu liên quan.

PostgreSQL \cite{postgresql} là hệ quản trị cơ sở dữ liệu quan hệ mã nguồn mở có độ ổn định cao, tuân thủ chặt chẽ các thuộc tính ACID, đảm bảo tính nhất quán và toàn vẹn dữ liệu.

Để tương tác với cơ sở dữ liệu, hệ thống sử dụng Sequelize ORM \cite{sequelize}. Sequelize giúp chuyển đổi các thao tác cơ sở dữ liệu thành các lời gọi hàm JavaScript, hỗ trợ quản lý mô hình dữ liệu, migration và góp phần hạn chế nguy cơ tấn công SQL Injection.
\subsection{MinIO trong lưu trữ tệp tin}
\label{subsection:3.2.3}
Để quản lý các tệp tin do người dùng tải lên, hệ thống sử dụng MinIO như một giải pháp lưu trữ đối tượng (Object Storage).

MinIO tương thích với chuẩn Amazon S3, cho phép lưu trữ và quản lý tệp tin một cách hiệu quả. Việc tách biệt dữ liệu tệp tin khỏi cơ sở dữ liệu giúp giảm tải cho máy chủ và tăng khả năng mở rộng của hệ thống.
\subsection{JWT và Swagger}
\label{subsection:3.2.4}
Hệ thống sử dụng JSON Web Token (JWT) \cite{jwt} để xây dựng cơ chế xác thực và phân quyền người dùng.

Sau khi đăng nhập thành công, người dùng được cấp một token để xác thực các yêu cầu tiếp theo mà không cần duy trì phiên làm việc (session) trên máy chủ. Ngoài ra, để đảm bảo tính an toàn cho tài khoản người dùng, mật khẩu khi lưu trữ vào cơ sở dữ liệu cũng được băm mã hóa một chiều bằng thư viện Bcrypt \cite{bcrypt}.

Bên cạnh đó, hệ thống sử dụng Swagger để xây dựng tài liệu API. Thư viện `swagger-jsdoc` và `swagger-ui-express` giúp tự động sinh tài liệu và cung cấp giao diện trực quan để kiểm thử các API ngay trên trình duyệt.

\end{document}